\chapter{Introducción}

% Presenta el problema, su contexto y el alcance de la tesis.
\section{Contexto histórico y científico}
[Desarrollar el contexto histórico y científico del trabajo.]

\section{Problema de investigación}
[Describir con precisión el problema que aborda la tesis.]

\section{Objetivos}
\subsection{Objetivo general}
Establecer un marco teórico que integre la noción de causalidad de Rosen y la jerarquía causal de Pearl para determinar si los grandes modelos de lenguaje tienen la capacidad de realizar inferencias causales legítimas.

\subsection{Objetivos específicos}
\begin{enumerate}[align=left]
  \item [Objetivo específico 1]
  \item [Objetivo específico 2]
  \item [Objetivo específico 3]
\end{enumerate}

\section{Preguntas de investigación}
\subsection{Pregunta central}
¿Qué criterios permiten distinguir la inferencia causal genuina de una asociación estadística simulada en el lenguaje?

\subsection{Preguntas secundarias}
\begin{enumerate}[align=left]
  \item [Pregunta secundaria 1]
  \item [Pregunta secundaria 2]
\end{enumerate}

\section{Hipótesis}
[Formular la hipótesis central y, si corresponde, las hipótesis secundarias.]

\section{Estructura de la tesis}
[Explicar brevemente el contenido de cada capítulo.]
